%Ajustes na legenda da figura. Incluindo espacamento apos a legenda
\captionsetup[figure]{labelformat={default},labelsep=period,font=footnotesize, name=\footnotesize{Fig.},justification=centering,singlelinecheck=false,belowskip=0\normalbaselineskip}
%\pagenumbering{gobble}

%Ajustes na legenda da tabela. 
%\captionsetup[table]{labelformat={default},labelsep=newline,font={sc,footnotesize},justification=centering,singlelinecheck=false}
\captionsetup[table]{format=plain,labelformat=simple,labelsep=period,justification=raggedright,singlelinecheck=false,skip=0pt,font=small}

%\renewcommand{\headrulewidth}{0pt}

\makeatletter
\newcommand{\linebreakand}{%
	\baselineskip
	\end{@IEEEauthorhalign}
	\hfill\mbox{}\par
	\mbox{}\hfill\begin{@IEEEauthorhalign}
}
\makeatother


%%%%%%%%%%%%%%%%%%%%%%%%%%%%%%%%%%%%%%%%%%%%%%%%%%%%%%%%%%%%%%%%%%%%%%%%%
%    Configuracaoes de Idioma
%    Considerando babel = brazil
%%%%%%%%%%%%%%%%%%%%%%%%%%%%%%%%%%%%%%%%%%%%%%%%%%%%%%%%%%%%%%%%%%%%%%%%% 

\addto\captionsbrazil{
  \renewcommand{\abstractname}{Abstract}
  \renewcommand{\figurename}{Figura}
  \renewcommand{\tablename}{Tabela}
}
%
% %considerando babel = english
% \addto\captionsenglish{
%   \renewcommand{\tablename}{Table}
% }
\renewcommand{\lstlistingname}{Código}
\renewcommand{\lstlistlistingname}{Lista de Códigos}

% Configure listings style for C code
\lstdefinestyle{prog}{
	backgroundcolor=\color{white},
	basicstyle=\ttfamily\footnotesize,
	keywordstyle=\color{purple}\bfseries,
	stringstyle=\color{green!50!black},
	commentstyle=\color{gray}\itshape,
	emph={
		size_t,
		Vector2
	},
	emphstyle=\color{purple}\bfseries,
	emph={[2]
		Poly6Kernel,
		GetNeighborParticles
	},
	emphstyle={[2]\color{teal}\bfseries},
	tabsize=2,
	breaklines=true
}
